% Figure: Prandtl-Meyer function nu(M) in degrees versus Mach number,
% gamma=1.4. nu is the angle through which a sonic flow must turn to
% reach Mach M by isentropic expansion. Saturates at nu_max ~ 130.5 deg
% as M -> infinity.
\begin{figure}[htbp]
\centering
\begin{tikzpicture}
\begin{axis}[
  width=12cm, height=8cm,
  xlabel={Mach number $M$}, ylabel={Prandtl-Meyer angle $\nu(M)$ (deg)},
  xmin=1, xmax=10, ymin=0, ymax=135,
  axis lines=left,
  grid=major, grid style={enggray!20},
  tick label style={font=\scriptsize},
  label style={font=\small},
  legend style={font=\scriptsize, at={(0.97,0.03)}, anchor=south east},
  legend cell align=left,
  domain=1:10, samples=300,
]
% nu(M) = sqrt(6) * atan( sqrt( (M^2-1)/6 ) ) - atan( sqrt(M^2-1) )
% with sqrt((g+1)/(g-1)) = sqrt(2.4/0.4) = sqrt(6) for gamma=1.4.
% pgf atan returns degrees; sqrt(6)=2.449 multiplies the first term,
% which is already in degrees. The second atan is also in degrees.
\addplot[engblue, very thick] {
  2.449*atan( sqrt( (x^2 - 1)/6 ) ) - atan( sqrt(x^2 - 1) ) };
\addlegendentry{$\nu(M)$, $\gamma=1.4$}
% asymptote nu_max = 90*(sqrt(6)-1) = 130.45 deg
\addplot[engred, dashed, domain=1:10, samples=2] {130.45};
\addlegendentry{$\nu_{\max}=130.5^{\circ}$}
% reference point M=2, nu~26.4
\node[engblack, font=\scriptsize, anchor=west] at (axis cs:2.1,26.4)
  {$\nu(2)\approx26.4^{\circ}$};
\addplot[only marks, mark=*, engblack, mark size=1.4pt] coordinates {(2,26.38)};
\end{axis}
\end{tikzpicture}
\caption[The Prandtl-Meyer function]{The Prandtl-Meyer function
$\nu(M)$ for $\gamma=1.4$, the angle through which a sonic stream
($M=1$, $\nu=0$) must turn around a convex corner to reach Mach $M$ by
isentropic expansion. A flow already at $M_1$ that turns through a
further convex angle $\Delta\theta$ reaches $M_2$ from $\nu(M_2) =
\nu(M_1) + \Delta\theta$, so the curve is read as a turning budget
rather than a function of position. The expansion is isentropic, so the
stagnation pressure is conserved across the fan, the property that
distinguishes an expansion corner from a compression corner (which
shocks). The angle saturates at $\nu_{\max} = 90^{\circ}(\sqrt{6}-1)
\approx 130.5^{\circ}$ as $M\to\infty$: a sonic flow cannot turn
through more than this before reaching infinite Mach number and zero
pressure, the limit of expansion into vacuum.}
\label{fig:vol03ch13:prandtl-meyer}
\end{figure}
